\section{Discussion and Conclusion}
\label{sec:discussion}

% DRAFT by Claude (21 Oct 2026), written from the notes that were here and
% DECISIONS_LOG.md (#89, #96, #99, #102, #104-#108, #113). Brian: rewrite
% it in your own words, and disclose the AI assistance.
\todo{Draft by Claude: rewrite in your own words.}

\textbf{Why tags do not transfer.} Not because tag values mean different
things in different places: within the same hike, harder sac\_scale grades
are slower in every region where they are used (mountain\_hiking vs
hiking $-5\%$, 10 of 10 regions; harder grades $-11\%$, 8 of 8). What
varies is where tags exist: sac\_scale is on 1--96\% of walked trail,
depending on the region, and untagged trail is faster than tagged trail in
some regions and slower in others ($-3.5\%$ to $+5.1\%$ within a hike). So
tags inform only part of each region's trails, and boosted trees on tags,
which can fit region-specific tag combinations, lose the most in transfer
($+1.5\% \to +0.4\%$). Terrain exists on every trail, and its effect has
the same sign in all 17 regions (steppiest vs smoothest fifth, $-12\%$).

\textbf{What terrain captures.} Uneven grade within the segment, such as
rock steps and short pitches, plus difficulty beyond grade. Its gain is
largest on descents and level trail ($+2.3$ to $+2.8\%$, additive) and
smallest on steep climbs ($+0.8\%$), consistent with footing setting the
pace going down and effort going up. The signal is the trail's elevation
profile, not micro-relief: 1~m lidar adds little over the 5~m grid, and
STRIDE's off-path roughness and understory terms add nothing on trails or
make it worse ($-18\% \to -35\%$).

\textbf{Use.} Eq.~(\ref{eq:terrain}) needs only a DEM. Fitted on the other
16 regions, it gives a pace for all 3,369~km of White Mountains trail,
77\% of which no trace in our data covers (median 15.2~min/km, vs 12.8 for
Tobler). It travels less far than the full terrain model, though
($+0.6\%$ with a whole zone held out, $-0.3\%$ between East and West).

\textbf{Limits.} Our hikers are OSM uploaders, not a random sample. Sparse
traces read slow (Sec.~\ref{sec:data}); hiking functions depend on the
sampling resolution~\cite{goodwin2025application}; distance along the
mapped trail would fix the cut corners and raise speeds by about 1--3\%,
and the comparisons hold at every recording rate. Per-segment gains are a
few percent because GPS noise and the hiker dominate.
Weather~\cite{thornton2022daymet} (a factor in~\cite{huang2026modeling})
slows our hikers on hot days ($-3\%$
above 25$^\circ$C) and with snow cover ($-4\%$), but neither weather
($-0.0\%$) nor date and time of day improves predictions in new regions,
and once a hiker's pace on the day is known, weather adds nothing.

\textbf{Conclusion.} Across 17 North American regions, terrain from
elevation data predicts hiking speed in regions a model has never seen;
OSM tags, vegetation and satellite embeddings mostly help where it was
trained. Tobler's function is too fast on gentle trail and overstates the
speed-up downhill. A Tobler-style function with one DEM-derived input
keeps part of the terrain gain and needs only a DEM, though a hiker's own
pace matters most.
\todo{Data release (segments without user names) if the mentors agree.}
